\section{边界与延伸：哪些是原文，哪些是后话}

\subsection{原文留下的空白}

读完前七讲，读者可能会觉得一切都很顺：作用量、相空间、量子化、共形块、计算。
但原文自己也清楚，有几处它只给了纲领而没有证明。把它们列出来，比记住结论更重要。

\textbf{第一，函子性。} 原文在讨论缝合时只写了一句 "It is obvious at this point"，
意思是：把曲面沿边界粘起来，对应的 Hilbert 空间应当做张量积。这个陈述把理论提升为
一个“模函子”（modular functor），但原文没有给出证明。严格化是后来的事（Segal 的框架、
以及大量关于共形块的手术公式的工作）。

\textbf{第二，标架反常。} (7.6) 的相位在一般手术下如何追踪，原文只处理了简单情形，
把一般情形留给了后续。这一项后来成为 Reshetikhin--Turaev 不变量里“缎带结构”的来源。

\textbf{第三，半形式修正。} 第四讲说过，$k\to k+h^\vee$ 的平移原文没有在正文里做；
它是后世（以及物理上的标架反常）补上的。

\textbf{第四，量子群。} 原文在讨论穿刺点约束时提到“量子群”，但没有构造。
后来的发展表明，可积表示范畴与量子群 $U_q(\mathfrak g)$ 在 $q=e^{2\pi\ii/(k+h^\vee)}$
处的表示范畴密切相关；这是 Kazhdan--Lusztig 等人工作的主题。

\textbf{第五，一般群。} 原文主要处理 $SU(N)$ 与 $SO(N)$；一般单群的情形类似，
但归一化与可积性条件需要逐群检查。

\subsection{两个常见的误读}

\subsubsection{CS/WZW 不是全息对偶}

第一个误读是把 CS/WZW 的对应说成“三维引力与二维 CFT 的全息对偶”，并由此联想到
AdS$_3$/CFT$_2$。两者不是一回事。

CS/WZW 是\textbf{态空间的等同}：三维理论在曲面 $\Sigma$ 上的 Hilbert 空间，
与二维理论在 $\Sigma$ 上的共形块空间，是同一个向量空间。这里没有“体”与“边界”的
动力学关系，没有引力子，没有全息字典。三维 Chern--Simons 理论的作用量 (2.1) 不含度规，
它根本没有引力自由度。

AdS$_3$/CFT$_2$ 是\textbf{动力学等价}：三维反德西特空间里的量子引力，与它边界上的
二维共形场论，描述同一套物理。它有引力子、有 BTZ 黑洞、有边界条件与全息字典。
两者唯一的接触点是：三维引力在适当条件下可以写成 Chern--Simons 理论
（规范群取非紧的 $SL(2,\RR)\times SL(2,\RR)$ 或其推广）。但这是一个“改写”，不是“等同”：
它要求非紧群、要求特定的边界条件，且只覆盖引力的经典/相空间层面。

读者若在文献里看到“CS 是全息的”，要追问：是哪一个 CS、哪一个群、哪一层级的对应。

\subsubsection{Khovanov 同调不是本文的预言}

第二个误读是把 Jones 多项式的范畴化（Khovanov 同调）说成“Witten 理论的自然推论”。
也不是。

Khovanov 同调是把 Jones 多项式从一个数提升为一个复形（或一个分次向量空间），
使得取 Euler 特征（带符号的维数）时退回 Jones 多项式。它满足函子性：链环的协边
给出复形之间的映射。这些性质说明它是一个“范畴化”的对象，比多项式携带更多信息。

但 1989 年的文章里没有任何关于同调、链复形或范畴化的暗示。把 Chern--Simons 理论
与 Khovanov 同调联系起来，要等到 Witten 2011 年（把 Khovanov 同调解释为
“高一维理论的边界条件”）以及 Kapustin--Witten 2007 年（4d 扭曲超对称规范理论的框架）。
即使在那之后，“CS 的 Wilson 线给出 Khovanov 复形”仍然不是一个已证明的定理；
它是一族猜想的统称。

\subsection{仍然开放的问题}

\begin{enumerate}[leftmargin=*]
  \item \textbf{函子性的严格化。} 把“缝合对应张量积”变成定理，需要什么条件？
        对一般的 $G$ 与一般的曲面，模函子的构造在什么范围内已知？
  \item \textbf{一般手术下的标架追踪。} (7.6) 在任意手术下的完整形式，
        以及它与缎带结构的关系。
  \item \textbf{半形式的系统处理。} $k\to k+h^\vee$ 在几何量子化框架里的一致推导，
        包括奇点处的相交上同调。
  \item \textbf{量子群与共形块的精确对应。} 可积表示范畴与 $U_q(\mathfrak g)$ 表示范畴
        在单位根处的等价，其精确陈述与证明。
  \item \textbf{与 Khovanov 同调的关系。} 是否存在一个“范畴化的 Chern--Simons 理论”，
        其边界给出 Khovanov 复形？
\end{enumerate}

\subsection{一条时间线}

把本文放进一条时间线里，能看清哪些是“当时的问题”、哪些是“后来的补课”。

\begin{center}
\small
\begin{tabular}{@{}p{1.9cm}p{7.4cm}p{4.6cm}@{}}
\toprule
年份 & 事件 & 与本文的关系 \\
\midrule
1984 & Jones 发现新的纽结不变量（von Neumann 代数） & 本文要解释的对象 \\
1985 & Jones 的 Bulletin AMS 短文；Kauffman 的状态和模型 & 组合定义成熟 \\
1986–87 & Drinfeld–Jimbo 量子群；Turaev–Reshetikhin 框架 & 提供二维/代数侧的语言 \\
1988 & Atiyah 在 Weyl 讨论会提两问；Witten 完成本文 & 本文的直接动机 \\
1989 & Moore–Seiberg、TUY 证明 Verlinde 公式（数学侧） & 补上本文只断言的部分 \\
1991 & Reshetikhin–Turaev 不变量；ADPW 的模空间量子化 & 严格化与几何化 \\
1994 & Faltings 证明 Verlinde 公式 & 数学证明的另一条路 \\
2000 & Khovanov 同调 & 范畴化，超出本文范围 \\
2007/2011 & Kapustin–Witten；Witten 的范畴化文章 & 范畴化的物理来源 \\
\bottomrule
\end{tabular}
\end{center}

\noindent 从这条线可以看出：本文处在“物理给出构造、数学随后补证明”的中间位置。
它不是一个封闭的定理，而是一个纲领；纲领的价值在于它指出了正确的对象
（曲面上的共形块），而不在于它把每一步都证完了。

\subsection{文献地图}

读到这里，读者可能需要知道去哪里查细节。下面按主题列出最常用的几份资料；
它们覆盖的范围、符号约定与严格程度各不相同，配合使用最有效。

\begin{center}
\small
\begin{tabular}{@{}p{3.0cm}p{5.7cm}p{5.4cm}@{}}
\toprule
资料 & 适合用来查什么 & 注意 \\
\midrule
Witten 1989（本文） & 原始论证、公式编号、物理动机 & 符号约定与教材不同；半形式略去 \\
河野《场论与拓扑学》 & 仿射代数、KZ、手术公式、Verlinde & CS 归一化差 $2\pi k$ \\
Schottenloher & Verlinde 的数学地位与证明史 & 第 11 章；物理推导不重复 \\
Moore TASI 2019 & 相空间量子化、标架反常、RCFT & 只查需要的节，不必通读 \\
Atiyah--Bott 1982 & 模空间的辛结构与 Kähler 结构 & 经典文献，符号较老 \\
TUY 1989 / Faltings 1994 & Verlinde 的数学证明 & 难度高，先读综述 \\
\bottomrule
\end{tabular}
\end{center}

\noindent 三条使用纪律：比较 $k$ 之前先对齐 $\Tr$ 与作用量的归一化；
看到 $q=e^{2\pi\ii/(k+h^\vee)}$ 时确认对方是否已做半形式平移；
凡涉及“标架”“相位”，先分清是配分函数的相位还是 Wilson 线的相位。

\subsection{开放问题的现状}

前面列的五个问题，到今天各有不同程度的进展，但没有一个被完全解决。
函子性方面，Segal 的框架与模函子理论给出了严格定义，但对一般 $G$ 与一般曲面，
构造仍依赖技术假设；手术公式的严格化（Reshetikhin--Turaev）覆盖了闭流形与链环，
但“缝合 = 张量积”的完整陈述仍不是一句话能证完的。
标架与半形式方面，两者的等价性在物理文献里是标准做法，
但一般奇点情形下的统一处理仍需要相交上同调与抛物丛的技术。
量子群与共形块的对应，在单位根处有 Kazhdan--Lusztig 等人的深刻结果，
但“可积表示范畴 $\cong$ $U_q(\mathfrak g)$ 的某子范畴”的精确陈述依赖 $q$ 的选取与范畴的完备化，
文献里的版本并不完全一致。范畴化方向是过去二十年最活跃的：
Khovanov 同调与 Khovanov--Rozansky 同调，以及它们与四维超对称规范理论的联系，
已经构成独立的研究领域；本讲义只讲到“它是什么、与本文什么关系”。

\subsection{怎么读这篇文章}

最后给三条读法，供不同背景的读者选择。

\textbf{快速读法（一天）。} 读前言、第一讲、第二讲、第七讲的小结。抓住四个词：
作用量、相空间、共形块、融合规则。

\textbf{标准读法（一周）。} 第一讲到第七讲按顺序读，每讲的公式自己推一遍。
第三、四讲是技术核心；第七讲是计算核心。

\textbf{研究读法（数周）。} 以原文为主，配合河野俊丈《场论与拓扑学》与 Schottenloher
《共形场论的数学导引》；遇到符号约定不一致时，回到原文页图核对。重点核对三件事：
$k$ 的归一化、$\Tr$ 的表示、半形式修正是否算进 $k$。

\subsection{结语}

回看全书，主线只有一条：一个不依赖度规的三维作用量，在曲面上量子化之后，
给出的态空间恰好是二维共形场论的共形块；纽结是这条主线上的标记点，
而 Jones 多项式是这条主线的可计算影子。

这条主线之所以值得写成一本书，是因为它同时是三个领域的交汇：
拓扑（三维流形的分类）、代数（量子群与融合范畴）、分析（几何量子化与模空间）。
任何一个领域单独看，都能看到一部分；只有把三者放在一起，才能理解为什么 Jones 多项式
会“出现”在三维理论里。

如果读者读完只记住一句话，我希望是这一句：\textbf{Chern--Simons 理论把“三维的拓扑”
编码成“二维的共形数据”，而量子化的曲面正是这两者之间的翻译词典。}

\begin{quote}
\small
\textbf{本讲要点。} 原文的五处空白（函子性、标架、半形式、量子群、一般群）是后续工作的入口；
CS/WZW 是态空间等同，不是全息对偶；Khovanov 同调是范畴化，不是本文的预言。
\end{quote}
